% =====================================================================
%  frontmatter/abstract.tex — executive summary (Stage-3 editor)
%  Rules: 250–350 words; problem -> question/hypotheses -> method ->
%  key findings (with numbers) -> implication; then 5–7 keywords.
% =====================================================================
\chapter*{תקציר מנהלים}
\addcontentsline{toc}{chapter}{תקציר מנהלים}
\markboth{תקציר מנהלים}{}

בין \en{2022} ל־\en{2023} העלה בנק ישראל את הריבית מ־\pct{0.10} ל־\pct{4.75}, והיקף העסקאות בשוק הדירות ירד. במקום להוריד מחירים באופן גלוי אימצו היזמים מבצעי מימון, ובהם \deal{20/80}, \deal{20/50/30} והלוואות בלון שהריבית עליהן ממומנת בידי היזם. מבצעים אלה דוחים את מבחן יכולת ההחזר של הרוכש עד יום המסירה. בספטמבר \en{2026} עמד מלאי הדירות החדשות הלא־מכורות על \en{85,080}, ומועדי המסירה של העסקאות שנחתמו בשנים \en{2023–2025} מתרכזים ב״קיר מסירות״ בשנים \en{2026–2027}.

הדוח בוחן באיזו מידה מבנים אלה משחזרים את מנגנוני הכשל של משבר הסאב־פריים בארצות הברית ושל משבר הנדל״ן בספרד, ומהי עמידות היזמים והבנקים בישראל בעת קיר המסירות. נבחנות שלוש השערות: אשראי צל המסווה ירידת מחירים (\en{H1}), לחץ החזר של רוכשים ביום המסירה (\en{H2}), והסטת הסיכון המערכתי מהבנקים אל היזמים (\en{H3}).

הגישה משלבת סקירת ספרות נושאית, השוואה רב־פרמטרית בין שלושת המשקים, וסימולציה מבנית מכוילת שנבחרה במטריצת החלטה משוקללת. הסימולציה כוללת מודל ערך נוכחי נקי, סימולציית מונטה קרלו של יחס ההחזר, רשת מבחני קיצון לפרויקט מייצג ומודל שחיקת הון (\en{CET1}) לחמש הקבוצות הבנקאיות. עובדות מוצגות עם מקור, ותוצאות המודלים מוצגות כתרחיש.

\en{H1} נתמכת חלקית: בתרחיש, עסקת \deal{20/80} גלומה בה הנחה של \en{10.9\%–12.8\%}, אך ברמת מדד המחירים הפער מצטמצם לכ־\en{1.4\%–2.0\%}. \en{H2} נתמכת בתרחיש הבסיס: כ־\pct{36} מהרוכשים חוצים יחס החזר של \pct{40}, והתוצאה רגישה לריבית. \en{H3} אינה נתמכת לפי סף ההכרעה שנקבע מראש: בתרחיש החמור יחס כיסוי החוב של היזם יורד ל־\en{0.77}, אך ההפסד הבנקאי הוא \pct{21.3} מההון העצמי של המערכת בגרסה השמרנית (הסופרת את ירידת המחיר פעמיים) ו־\pct{14.2} בגרסה המבנית, מעל הסף של \pct{10}, הנבחן לפני מס. הקביעה שברירית: אחרי מס הגרסאות חלוקות (\pct{13.9} ו־\pct{9.2}) והתוצאה אינה מכריעה; בשיעור העברה של \en{0.5}, או בכושר ספיגה על מלוא ערך הפרויקט (הגדרת הפועלים), שתיהן נמוכות מ־\pct{10}; וכשל רוכשים לבדו אינו יוצר הפסד בנקאי. בתרחיש החמור יורדים עד ארבעה בנקים אל מתחת ליחס ההון המזערי, אך חריגה זו אינה חדלות פירעון, ובאף תרחיש שנבחן בפרמטרי הבסיס (עד \en{40\%/40\%/24}) אין בנק היורד אל מתחת ל־\pct{6.5}. ישראל דומה אפוא לספרד יותר מאשר לארצות הברית, והבנקים אינם מנותקים מהפסדי היזמים.

בשלב המימוש ייאספו נתוני מסירות, ביטולים ומשכנתאות ברמת הפרויקט, המודלים יכוילו מחדש מולם, ואומדן ההפסד הבנקאי יחודד בהתאם.

\vspace{1em}
\noindent\textbf{מילות מפתח:} מבצעי מימון קבלניים; קיר המסירות; אשראי צל; מבחני קיצון; חוק המכר; משבר הסאב־פריים; מלאי דירות לא־מכורות
